% !TEX root = cedar2020smartmeters.tex
% FICTIONAL TEST MATERIAL: this paper does not exist; it was written for the
% citation-audit fixture.
\documentclass{article}
\usepackage{amsmath}

\title{Inlet Smart Meters for Night-Flow Anomaly Detection}
\author{E. Fixture}
\date{Journal of Fixture Studies, 2020}

\begin{document}
\maketitle

\noindent\textbf{FICTIONAL TEST MATERIAL: this paper does not exist; it was
written for the citation-audit fixture.} No author, utility, venue, or
measurement described here is real.

\begin{abstract}
A new leak in a district adds a steady flow that is easiest to see at night,
when legitimate demand is lowest. We installed smart meters at the inlets of
ten districts of a fictional network and opened a hydrant in each to simulate
a new leak. A simple night-flow anomaly rule flagged the simulated leak within
24 hours in 9 of the 10 trials. Detecting a new leak within a day, rather than
at the next scheduled survey, shortens the time it runs before repair and so
reduces the water it loses.
\end{abstract}

\section{Introduction}
\label{sec:introduction}

Most leaks are not reported by the public. They run until an active survey
finds them, and surveys visit each district only every few months. The water a
leak loses grows with the time it runs, so detecting it early matters as much
as finding it precisely.

A district metered at its inlet offers a cheap signal. Between two and four in
the morning, almost all legitimate demand stops, and the inlet flow settles
near a stable minimum. A new leak raises that minimum. This paper tests
whether inlet smart meters can turn that signal into a timely alarm.

\section{Method}
\label{sec:method}

\subsection{Meters and districts}
\label{sec:method-meters}

Ten districts of the fictional Fixture Ridge network, each fed through a single
inlet, received a smart meter that logged flow every 15 minutes and reported it
daily. The districts ranged from 400 to 3,000 customer connections.

\subsection{Anomaly rule}
\label{sec:method-rule}

For each district, the rule computes the minimum night flow over a 30-day
rolling window. It raises an anomaly when the minimum night flow on a given
night exceeds the rolling median by more than three times the rolling
median absolute deviation.

\subsection{Simulated leaks}
\label{sec:method-trials}

In each district, a crew opened a hydrant at a fixed, small flow at an
unannounced time to simulate a new leak, and left it open for three days. The
analysts running the rule did not know when or where the hydrant was opened.
One trial was run per district, ten trials in total.

\section{Results}
\label{sec:results}

\subsection{Detection within 24 hours}
\label{sec:results-detection}

Smart meters at the district inlets flagged the night-flow anomaly from the
new leak within 24 hours in 9 of the 10 trials. In each of these nine trials
the alarm came on the first night after the hydrant was opened. In the tenth
trial, the largest district, the simulated leak was small compared with the
normal night-flow variation and the rule flagged it only on the third night.

\subsection{False alarms}
\label{sec:results-false}

Over the 30 days before the trials, the rule raised two false alarms across all
ten districts, both on nights with large, known customer demands such as a
swimming-pool fill.

\subsection{Effect on run time}
\label{sec:results-runtime}

Without the meters, a new leak in these districts would run until the next
active survey, on average several weeks later. Detecting it within a day
shortens its run time by the same amount and reduces the water it loses before
repair accordingly.

\section{Limitations}
\label{sec:limitations}

This study tests detection only, with simulated leaks.

This study does not examine pressure control: district pressures were left
unchanged, and we make no claim about how pressure management affects leakage.
This study does not examine bursts: a burst is obvious without a meter, and the
trials simulated only small new leaks. The meters say that a district has a new
leak but not where it is; locating it still needs a field survey. Ten trials
are enough to show the rule works but not to estimate its detection rate
precisely.

\end{document}
